\section{Spectroscopy step by step}
\label{sec:spec-guide}

\subsection{Slit spectroscopy}

\begin{enumerate}
  \item Mode \code{slit}; set the resolving power $R$ (or load a measured
        slit-width/$R$ curve, which then overrides it), the slit width and
        the extraction height in arcsec, and pixels per resolution element.
  \item Set the wavelength range and the S/N reference wavelength: the
        time-series evaluates that wavelength bin at every time sample.
  \item Choose the \textbf{slit orientation}. \code{parallactic} assumes
        you rotate the slit to the parallactic angle (reported in the
        Results window for every time), avoiding atmospheric-dispersion
        losses. \code{fixed} applies the worst-case per-wavelength loss ---
        expect tens of per cent missing at the blue end at airmass 1.5--2.
  \item Run. The plot window shows the detected source rate and S/N per
        resolution element, with a manual time slider through the night.
\end{enumerate}

\subsection{Slitless spectroscopy and the Star Analyser}

Mode \code{slitless}. Two ways to describe the disperser:

\begin{itemize}
  \item \textbf{Star Analyser 100/200 (recommended)}: enter the groove
        density (100 or 200 lines/mm) and the grating-to-sensor distance in
        mm; the dispersion in \si{\angstrom}/pixel follows from the
        physics (an SA100 at \SI{42}{mm} on \SI{4.63}{\micro\metre} pixels
        gives $\simeq11$ \si{\angstrom}/pixel). Optionally set the
        first-order grating efficiency ($\sim0.5$ is realistic).
  \item \textbf{Manual}: leave the groove density at 0 and type the
        dispersion in \si{\angstrom}/pixel directly.
\end{itemize}

The resolution element is set by seeing, the intrinsic grating LSF and
atmospheric dispersion; with $2$--$3''$ seeing expect $R\sim100$--200
regardless of exposure. The cross-dispersion extraction height is the only
aperture; there are no slit losses. Slitless results are flagged
\emph{experimental} until validated on a calibrated instrument.

\subsection{Reading the spectroscopic output}

All spectroscopic counts and S/N are \textbf{per resolution element}, not
per pixel. The Results table lists, per wavelength: the element width, the
effective $R$, source and sky rates, S/N, peak electrons and ADU of the
brightest pixel, the saturation flag and the maximum unsaturated exposure.
